% =====================================================================
%  Section: Network-driven excitation of the drivetrain torsional mode
%           by a LEOGO process-load event (Low-Emission scenario)
%  Follows on from \label{sec:forcedresponse}.
%  Required in preamble:
%    \usepackage{amsmath, amssymb}
%    \usepackage{graphicx}
%    \usepackage{booktabs}
%    \usepackage{siunitx}
%  Adjust the figure path (fig/...) to your own.
% =====================================================================

\section{Network-Driven Excitation from a Platform Process Load}
\label{sec:network-excitation}

The forced-response study of \S\ref{sec:forcedresponse} established that the
drivetrain torsional mode at \SI{3.49}{\hertz} can be driven to resonance by an
electrical disturbance applied \emph{at the turbine terminal}. That experiment
was a controlled probe: a clean single-frequency shunt modulation used to trace
the electromechanical transfer path and locate the resonance. It does not yet
answer the question that matters for an integrated offshore energy system such
as LEOGO: can a disturbance originating \emph{elsewhere in the platform grid}
--- a process-load event on the oil-and-gas side --- reach the wind turbine and
excite the same mode? This section answers that question in the affirmative,
using the reduced model in which the electromechanical path is open by
construction.

\subsection{Point of common coupling}
\label{subsec:ne-pcc}

On the LEOGO platform every large process load --- the gas-export compressor
trains (\texttt{GEX}, \SI{9.7}{\mega\watt} each), the drilling loads, the
utility and pump loads --- and the wind park are supplied from the same
\SI{11}{\kilo\volt} main switchboard, \emph{Main~Bus~A}. That busbar is
therefore the point of common coupling (PCC) between the process side and the
wind turbine. A power fluctuation produced by a process load appears at the PCC
and propagates, through the intervening network impedance, to the turbine
terminal. Accordingly the disturbance in this section is injected as a
shunt-load admittance at Main~Bus~A rather than at the turbine, so that it must
traverse the platform network before it can act on the drivetrain. Because the
individual constant-impedance process loads are folded into the reduced network,
injecting the aggregate disturbance at the PCC is the faithful, lumped
representation of a process-load event.

\subsection{Network-to-turbine transfer}
\label{subsec:ne-transfer}

Repeating the frequency sweep of \S\ref{subsec:fr-simplified} with the
disturbance moved from the turbine terminal to the PCC traces out the
network-to-turbine transfer of Figure~\ref{fig:ne-curve} and
Table~\ref{tab:ne-sweep}. The response remains sharply resonant and still peaks
\emph{exactly} at \SI{3.49}{\hertz}: the torsional mode is reached from the
LEOGO side and rings at its own frequency, independent of the injection point.
The intervening network impedance attenuates the amplitude --- at resonance the
peak-to-peak shaft torque is \SI{0.60}{\mega\newton\metre} from the PCC versus
\SI{7.96}{\mega\newton\metre} for the same disturbance applied directly at the
terminal, a factor of about $13$ --- but the shape and location of the resonance
are preserved (Figure~\ref{fig:ne-compare}). Normalising each curve to its own
peak (Figure~\ref{fig:ne-normalized}) makes the point explicit: the two curves
are the same resonance, merely scaled. The attenuation grows with frequency, as
expected for the series reactance of the feeders and step-up transformers
between the PCC and the turbine.

\begin{figure}[t]
  \centering
  % \includegraphics[width=0.8\linewidth]{fig/WT_LEOGO_network2turbine_curve.png}
  \caption{Network-to-turbine transfer: peak-to-peak shaft torque versus the
  frequency of an electrical disturbance injected at the PCC (Main~Bus~A). The
  response peaks at the \SI{3.49}{\hertz} torsional mode, reached from the LEOGO
  side.}
  \label{fig:ne-curve}
\end{figure}

\begin{table}[t]
  \centering
  \caption{Frequency sweep with the disturbance injected at the PCC
  (Main~Bus~A), compared with direct terminal injection
  (\S\ref{subsec:fr-simplified}). Both use $\hat{P}=\SI{2.0}{\mega\watt}$.}
  \label{tab:ne-sweep}
  \begin{tabular}{lccc}
    \toprule
    $f$ [\si{\hertz}] & PCC $\Delta T_{\mathrm{sh}}$
      [\si{\mega\newton\metre}] & Terminal $\Delta T_{\mathrm{sh}}$
      [\si{\mega\newton\metre}] & Attenuation \\
    \midrule
    1.00 & 0.051 & 0.22 & $4.4\times$ \\
    2.00 & 0.142 & 0.66 & $4.7\times$ \\
    2.50 & 0.191 & 1.23 & $6.4\times$ \\
    3.00 & 0.299 & 2.65 & $8.9\times$ \\
    3.25 & 0.420 & 4.46 & $10.6\times$ \\
    \textbf{3.49} & \textbf{0.595} & \textbf{7.96} & $\mathbf{13.4\times}$ \\
    3.75 & 0.326 & 5.53 & $17.0\times$ \\
    \bottomrule
  \end{tabular}
\end{table}

\begin{figure}[t]
  \centering
  % \includegraphics[width=0.8\linewidth]{fig/WT_LEOGO_network2turbine_compare.png}
  \caption{Resonance curves for direct terminal injection (left axis) and
  PCC injection (right axis). Both peak at \SI{3.49}{\hertz}; only the amplitude
  differs, owing to the network attenuation.}
  \label{fig:ne-compare}
\end{figure}

\begin{figure}[t]
  \centering
  % \includegraphics[width=0.8\linewidth]{fig/WT_LEOGO_network2turbine_normalized.png}
  \caption{The two curves of Figure~\ref{fig:ne-compare} normalised to their own
  peaks: identical resonance shape, confirming that the same torsional mode is
  excited whether the disturbance enters at the turbine or at the PCC.}
  \label{fig:ne-normalized}
\end{figure}

\subsection{A realistic process-load pulsation}
\label{subsec:ne-pulsation}

The swept sinusoid is a measurement tool, not a physical event. To make the
scenario concrete, the disturbance is now interpreted as the fluctuating
component of a cyclically loaded process machine --- a reciprocating or
periodically loaded compressor whose electrical power pulsates about its mean
draw. A modest pulsation amplitude of $\hat{P}=\SI{0.5}{\mega\watt}$ (a few
percent of a \SI{9.7}{\mega\watt} compressor train) at \SI{3.49}{\hertz} is
applied at the PCC. The drivetrain builds up to a sustained shaft-torque
oscillation of amplitude \SI{74}{\kilo\newton\metre}
(Figure~\ref{fig:ne-pulsation}), i.e.\ a transfer gain of
\begin{equation}
  \frac{\Delta \hat{T}_{\mathrm{sh}}}{\hat{P}}
  \approx \SI{149}{\kilo\newton\metre\per\mega\watt}
  \label{eq:ne-gain}
\end{equation}
of shaft-torque swing per megawatt of pulsating power at the PCC. Even a small,
entirely realistic load ripple therefore produces a clearly measurable cyclic
loading of the wind-turbine drivetrain. Following the disturbance along its path
(Figure~\ref{fig:ne-propagation}) shows the chain explicitly: the pulsation
modulates the Main~Bus~A voltage, which modulates the wind-turbine terminal
power, which through $T_e$ drives the shaft torque, all locked to the
\SI{3.49}{\hertz} mode.

\begin{figure}[t]
  \centering
  % \includegraphics[width=0.8\linewidth]{fig/WT_LEOGO_process_pulsation_response.png}
  \caption{A \SI{0.5}{\mega\watt} process-load pulsation at \SI{3.49}{\hertz}
  at the PCC (top) and the resonant build-up of the drivetrain shaft-torque
  deviation it drives (bottom).}
  \label{fig:ne-pulsation}
\end{figure}

\begin{figure}[t]
  \centering
  % \includegraphics[width=0.8\linewidth]{fig/WT_LEOGO_process_pulsation_propagation.png}
  \caption{Propagation of the pulsation from the PCC to the turbine (a few
  steady-state cycles): Main~Bus~A voltage, wind-turbine terminal power and
  shaft torque, all oscillating at \SI{3.49}{\hertz}.}
  \label{fig:ne-propagation}
\end{figure}

\subsection{Frequency selectivity of the coupling}
\label{subsec:ne-selectivity}

The transfer sweep of \S\ref{subsec:ne-transfer} already implies that the
network-to-drivetrain coupling is sharply frequency dependent, but this is worth
stating explicitly, because it clarifies the nature of the interaction. The
grid does not couple broadband energy into the drivetrain; the shaft responds
strongly only to disturbances whose spectral content coincides with the
\SI{3.49}{\hertz} eigenfrequency. Figure~\ref{fig:ne-selectivity} makes this
concrete by applying the \emph{same} \SI{0.5}{\mega\watt} PCC pulsation of
\S\ref{subsec:ne-pulsation} at three frequencies on a shared torque axis. Off
resonance the drivetrain barely reacts --- a steady swing of
\SI{16}{\kilo\newton\metre} at \SI{2.0}{\hertz} and \SI{29}{\kilo\newton\metre}
at \SI{3.0}{\hertz} --- whereas the on-resonance disturbance at
\SI{3.49}{\hertz} rings the mode up to \SI{74}{\kilo\newton\metre}, roughly
$4.8\times$ the \SI{2.0}{\hertz} amplitude for identical forcing.

This selectivity resolves an apparent tension with the modal analysis of
\S\ref{sec:frequency-analysis}, where the torsional mode is found to be
electromechanically \emph{isolated}: no system eigenmode spontaneously ties the
platform grid to the drivetrain. The two findings are consistent. Modal coupling
asks whether a single eigenvector spans both subsystems, and the answer is no.
The excitation studied here is instead a \emph{forced} transmission along the
open path $\Delta P \!\to\! \Delta T_e \!\to\! \Delta T_{\mathrm{sh}}$: a grid
disturbance drives the lightly damped mode only in proportion to the spectral
energy it carries at the modal frequency. Because the mode has a quality factor
of only $Q \approx 11$ ($\zeta \approx \SI{4.4}{\percent}$), its resonant window
is broad rather than needle-sharp --- hence the \SI{3.0}{\hertz} pulsation
already produces an appreciable response --- but the peak remains firmly centred
on \SI{3.49}{\hertz}. It is precisely this decoupling that makes the deliberate
frequency matching of \S\ref{subsec:ne-pulsation} necessary: the drivetrain is
excited because the disturbance is placed at its resonance, not because the two
systems are inherently coupled.

\begin{figure}[t]
  \centering
  % \includegraphics[width=0.8\linewidth]{fig/WT_LEOGO_process_frequency_selectivity.png}
  \caption{Frequency selectivity of the network-to-drivetrain coupling: the same
  \SI{0.5}{\mega\watt} PCC pulsation applied at \SI{2.0}{\hertz} (off resonance),
  \SI{3.0}{\hertz} (off resonance) and \SI{3.49}{\hertz} (on resonance), plotted
  on a common shaft-torque axis. Only the disturbance at the eigenfrequency rings
  the mode.}
  \label{fig:ne-selectivity}
\end{figure}

\subsection{A discrete load-switching event}
\label{subsec:ne-step}

Resonant excitation does not require a tuned, sustained pulsation: any
disturbance with spectral content near the eigenfrequency will ring a lightly
damped mode. To demonstrate this, a discrete load-switching event is applied at
the PCC --- a \SI{9.7}{\mega\watt} compressor block energised over a
\SI{20}{\milli\second} breaker time, modelled as a ramped shunt load,
\begin{equation}
  y_{\mathrm{sh}}(t) = y_{\mathrm{load}}\,
  \min\!\Big(\tfrac{t-t_{\mathrm{on}}}{t_{\mathrm{r}}},\,1\Big)\,
  \mathbb{1}[t \ge t_{\mathrm{on}}].
  \label{eq:ne-step}
\end{equation}
The step response, shown in Figure~\ref{fig:ne-step}, has two components. The
dominant, slow part is the electromechanical settling of the LEOGO synchronous
generators and their governors as they absorb the \SI{9.7}{\mega\watt} load
change over several seconds. Superimposed on it is a decaying oscillation: the
broadband step excites the torsional mode, which then rings. High-pass filtering
the response to remove the slow settling isolates this ring (lower panel): it
oscillates at \SI{3.47}{\hertz} with a peak-to-peak amplitude of
\SI{132}{\kilo\newton\metre} and decays with a fitted damping ratio of
\SI{4.3}{\percent}. That value matches the \SI{4.37}{\percent} obtained
independently from the eigenvalue analysis (\S\ref{sec:frequency-analysis}),
confirming that the ringing is genuinely the identified torsional mode and not
an artefact of the forcing. A single switching event on the process side thus
leaves a measurable torsional transient on the wind-turbine shaft.

\begin{figure}[t]
  \centering
  % \includegraphics[width=0.8\linewidth]{fig/WT_LEOGO_process_step_ringdown.png}
  \caption{A discrete \SI{9.7}{\mega\watt} load-switching event at the PCC. Top:
  the full shaft-torque deviation, dominated by the slow genset/governor
  settling. Bottom: the high-pass-isolated torsional ring at \SI{3.47}{\hertz}
  with its fitted decay envelope ($\zeta=\SI{4.3}{\percent}$), matching the modal
  damping.}
  \label{fig:ne-step}
\end{figure}

\subsection{Conclusion}
\label{subsec:ne-conclusion}

Taken together, the three experiments close the loop between the electrical
network and the wind-turbine drivetrain. The transfer sweep shows that the
\SI{3.49}{\hertz} torsional mode is reachable from the platform side of the grid,
attenuated but unchanged in frequency. A realistic \SI{0.5}{\mega\watt}
compressor pulsation at that frequency drives a \SI{74}{\kilo\newton\metre}
shaft-torque oscillation, and a single \SI{9.7}{\mega\watt} switching event rings
the mode with a decay rate that matches the eigenvalue prediction. For an
integrated low-emission energy system, in which offshore wind supplies an
oil-and-gas platform with large rotating process loads, this establishes a
concrete electromechanical interaction pathway: periodic or transient
disturbances on the process side can couple, through the shared network, into
cyclic mechanical loading of the wind-turbine drivetrain. Whether such an
interaction is benign or damaging depends on how close the platform's
characteristic load frequencies lie to the drivetrain torsional frequency ---
the quantitative screening of which is the natural continuation of this work.
